% Section 12.5: Accidental Symmetries.
% Source: physical PDF pp. 552--554 (printed pp. 529--531), ending before Problems.
% Equations: (12.5.1)--(12.5.8). Figures/tables: none.
% Footnotes: two, including the section-title note.

\subsection[Accidental Symmetries]{Accidental Symmetries\texorpdfstring{\protect\footnotemark}{}}
\footnotetext{This section lies somewhat out of the book's main line of
development, and may be omitted in a first reading.}

In Section 12.3 we saw that there are good reasons to adopt renormalizable
field theories as approximate descriptions of nature at sufficiently low
energy. It often happens that condition of renormalizability is so stringent
that the effective Lagrangian automatically obeys one or more symmetries,
which are not symmetries of the underlying theory, and may therefore be
violated by the suppressed non-renormalizable terms in the effective
Lagrangian. Indeed, most of the experimentally discovered symmetries of
elementary particle physics are \textquotedblleft accidental
symmetries\textquotedblright{} of this sort.

A classic example is provided by the inversions and flavor conservation in the
electrodynamics of charged leptons. The most general renormalizable and gauge-
and Lorentz-invariant Lagrangian density for photons and fields $\psi_i$ of
spin $\tfrac12$ and charge $-e$ takes the form
\begin{align}
\mathcal L={}&-\frac14 Z_3F_{\mu\nu}F^{\mu\nu}\notag\\
&-\sum_{ij}Z_{Lij}\bar\psi_{Li}[\sl{\partial}+ie\sl A]\psi_{Lj}
-\sum_{ij}Z_{Rij}\bar\psi_{Ri}[\sl{\partial}+ie\sl A]\psi_{Rj}\notag\\
&-\sum_{ij}M_{ij}\bar\psi_{Li}\psi_{Rj}
-\sum_{ij}M_{ij}^{\dagger}\bar\psi_{Ri}\psi_{Lj},
\tag{12.5.1}\label{eq:12.5.1}
\end{align}
where $i,j$ are summed over the three lepton flavors ($e$, $\mu$, and $\tau$),
$\psi_{iL}$ and $\psi_{iR}$ are the left- and right-handed parts of the field
$\psi_i$, defined by
\begin{equation}
\psi_{Li}=\frac12(\mathbf 1+\gamma_5)\psi_i,
\qquad
\psi_{Ri}=\frac12(\mathbf 1-\gamma_5)\psi_i,
\tag{12.5.2}\label{eq:12.5.2}
\end{equation}
and $Z_L$, $Z_R$, and $M$ are numerical matrices. We are not assuming
anything about lepton flavor conservation, so the matrices $Z_{Lij}$,
$Z_{Rij}$ and $M_{ij}$ need not be diagonal. Also we are not assuming
anything about invariance under P, C, or T invariance, so there is no necessary
relation between $Z_L$ and $Z_R$, or between $M$ and $M^\dagger$. The only
constraints on these matrices come from the reality of the Lagrangian density,
which requires that $Z_{Lij}$ and $Z_{Rij}$ are Hermitian, and from the
canonical anticommutation relations, which require that $Z_{Lij}$ and
$Z_{Rij}$ are positive-definite.

Now suppose we replace the lepton fields $\psi_L$, $\psi_R$ with new fields
$\psi_L'$, $\psi_R'$ defined by
\begin{equation}
\psi_L=S_L\psi_L',
\qquad
\psi_R=S_R\psi_R',
\tag{12.5.3}\label{eq:12.5.3}
\end{equation}
where $S_{L,R}$ are non-singular matrices that can be chosen as we like. The
Lagrangian density when expressed in terms of these new fields then takes the
same form as in Eq.~\eqref{eq:12.5.1}, but with new matrices
\begin{equation}
Z_L'=S_L^\dagger Z_LS_L,
\qquad
Z_R'=S_R^\dagger Z_RS_R,
\qquad
M'=S_L^\dagger MS_R.
\tag{12.5.4}\label{eq:12.5.4}
\end{equation}
We can choose $S_L$ and $S_R$ so that $Z_L'=Z_R'=\mathbf 1$. (Take
$S_{L,R}=U_{L,R}D_{L,R}$, where $U_{L,R}$ are the unitary matrices that
diagonalize the positive-definite Hermitian matrices $Z_{L,R}$, and the
$D_{L,R}$ are the diagonal matrices whose elements are the inverse square
roots of the eigenvalues of $Z_{L,R}$.)

Now make another transformation, to lepton fields $\psi_L''$ defined by
\begin{equation}
\psi_L'=S_L'\psi_L'',
\qquad
\psi_R'=S_R'\psi_R''.
\tag{12.5.5}\label{eq:12.5.5}
\end{equation}
The Lagrangian density again takes the same form when expressed in terms of
these new fields, with new matrices
\begin{equation}
Z_L''=S_L'{}^\dagger S_L',
\qquad
Z_R''=S_R'{}^\dagger S_R',
\qquad
M''=S_L'{}^\dagger M'S_R'.
\tag{12.5.6}\label{eq:12.5.6}
\end{equation}
This time we take $S_{L,R}'$ unitary, so that again
$Z_L''=Z_R''=\mathbf 1$. We choose these unitary matrices so that $M''$ is
real and diagonal. (By the polar decomposition principle, $M'$ like any
square matrix may be put in the form $M'=VH$, where $V$ is unitary and $H$ is
Hermitian. Take $S_L'=S_R'{}^\dagger V^\dagger$ and choose $S_R'$ as the
unitary matrix that diagonalizes $H$.) Dropping primes, the Lagrangian density
now takes the form
\begin{align}
\mathcal L={}&-\frac14 Z_3F_{\mu\nu}F^{\mu\nu}
-\sum_i\bar\psi_{Li}[\sl{\partial}+ie\sl A]\psi_{Li}
-\sum_i\bar\psi_{Ri}[\sl{\partial}+ie\sl A]\psi_{Ri}\notag\\
&-\sum_i m_i\bar\psi_{Li}\psi_{Ri}
-\sum_i m_i\bar\psi_{Ri}\psi_{Li},
\tag{12.5.7}\label{eq:12.5.7}
\end{align}
where $m_i$ are real numbers, the eigenvalues of the Hermitian matrix $H$.
Finally, this can be put in the more familiar form
\begin{equation}
\mathcal L=-\frac14Z_3F_{\mu\nu}F^{\mu\nu}
-\sum_i\bar\psi_i[\sl{\partial}+ie\sl A]\psi_i
-\sum_i m_i\bar\psi_i\psi_i.
\tag{12.5.8}\label{eq:12.5.8}
\end{equation}
With the Lagrangian taking this form, it is now apparent that any
renormalizable Lagrangian for lepton electrodynamics automatically conserves
P, C, and T, as well as the numbers of leptons (minus the numbers of
antileptons) of each flavor: electron, muon, and tauon.\footnote{This was first
shown by Feinberg, Kabir, and myself~\hyperref[ref:12.21]{[21]}.
Feinberg~\hyperref[ref:12.22]{[22]} had earlier noted that weak interaction
effects in a theory with only one neutrino species would give rise to an
observable rate for the process $\mu\to e+\gamma$, a difficulty that was only
resolved by the discovery of a second neutrino species.} In particular,
despite the appearance of Eq.~\eqref{eq:12.5.1}, this theory does not allow
such processes as $\mu\to e+\gamma$. The reader may perhaps worry whether it
is correct to identify the lepton fields as the $\psi_i$ (previously called
$\psi_i''$) appearing in Eq.~\eqref{eq:12.5.8}, which obviously conserves
lepton flavor, rather than the $\psi_i$ appearing in Eq.~\eqref{eq:12.5.1},
which seems to allow processes like $\mu\to e+\gamma$.

Such worries may be put aside; as stressed in Section 10.3, there is no one
field that can be identified as \emph{the} field of the electron or muon. In
fact, although Eq.~\eqref{eq:12.5.1} yields a non-vanishing matrix element for
radiative decay of lepton 1 into lepton 2 \emph{off} the lepton mass shell, by
taking the lepton momenta on the mass shell we find a vanishing $S$-matrix for
all such processes even when calculated using Eq.~\eqref{eq:12.5.1}.

It was essential in deriving these results that the same electric charge
appeared in Eq.~\eqref{eq:12.5.1} for both the left- and right-handed parts of
the lepton fields or, in other words, that both left- and right-handed parts of
the lepton fields transform in the same way under electromagnetic gauge
transformations. As we shall see in Volume II, for similar reasons the modern
renormalizable theory of strong interactions known as quantum chromodynamics
automatically conserves C, and (aside from certain non-perturbative effects) P
and T, as well as the numbers of quarks (minus the numbers of antiquarks) of
each quark flavor. We shall also see in Volume II that the simplest version of
the renormalizable standard model of weak and electromagnetic interactions
automatically conserves lepton flavor (though not C and P) for reasons similar
to those described here for electrodynamics. It remains an open possibility
that non-renormalizable interactions arising from higher mass scales may
violate any of these conservation laws.

